Una espira magnètica es troba situada en el pla $YZ$, té un radi $R = 5\ \mathrm{cm}$ i transporta un corrent de 10~A.

\figura[0.3]{fig1.png}

\begin{apartat}{1}
Calculeu el mòdul del camp magnètic en el centre de l'espira (en~$\mathrm{\mu T}$).
\end{apartat}

\begin{apartat}{1}
Quin sentit ha de tenir el corrent elèctric que circula per l'espira perquè el camp magnètic en el centre vagi en el sentit positiu de l'eix $x$?
\end{apartat}

\blocetiquetat{Dada:}{$\mu_0 = 4\pi\times 10^{-7}\ \mathrm{T\,m\,A^{-1}}$.}
\nota{El mòdul del camp magnètic creat per una espira magnètica en un punt de l'eix $x$ és:
\[ B(x) = \frac{\mu_0 I}{2}\,\frac{R^2}{(x^2+R^2)^{3/2}}. \]}
